% STARD Checklist for "Large language models for self-administered clinical vignette assessment"
% Daniels, Zhang, Nguyen, Duong

\documentclass{article}

\usepackage{booktabs}
\usepackage{tabularx}
\usepackage[margin=1in,footskip=0.25in]{geometry}
\usepackage{hyperref}
\usepackage{setspace}
  \doublespacing

\title{STARD Checklist: Large Language Models for Self-Administered Clinical Vignette Assessment}

\author{
  Benjamin Daniels$^{1,2,\ast}$,
  Wenjia(Stella) Zhang$^{1}$,
  Hoang Nguyen$^{3,4}$,
  David Duong$^{5,6}$ \\[6pt]
  \small $^{1}$ Harvard T.H.\ Chan School of Public Health, Boston, MA, USA \\
  \small $^{2}$ gui2de, Georgetown University, Washington, DC, USA \\
  \small $^{3}$ Partnership for Health Advancement in Vietnam (HAIVN), Hanoi, Vietnam \\
  \small $^{4}$ Beth Israel Deaconess Medical Center (BIDMC), Boston, MA, USA \\
  \small $^{5}$ Harvard Medical School, Boston, MA, USA \\
  \small $^{6}$ Brigham and Women's Hospital, Boston, MA, USA \\[4pt]
  \small $^{\ast}$ Corresponding author: bdaniels@fas.harvard.edu
}

\date{}

\begin{document}

\maketitle

\section*{STARD 2015 Checklist}

This document presents a completed STARD (Standards for Reporting Diagnostic Accuracy Studies) checklist for the diagnostic accuracy validation component of the study, in which LLM-automated coding of clinical vignette transcripts (the index test) was compared against human-coded ground truth (the reference standard). The diagnostic question evaluated is: ``Can an LLM-automated system accurately code essential diagnostic checklist items from clinical vignette transcripts?''

\subsection*{Index Test Description}

Index test: Anthropic Claude language model (claude-haiku-4-5) automated coding of pre-specified binary checklist items (whether the provider asked each history question, ordered each recommended examination/test, or made the correct diagnosis) from interaction transcripts.

\subsection*{Reference Standard}

Reference standard: Manual human coding of the same checklist items from the same transcripts by the research team, using an identical classification schema.

\subsection*{Patient Population}

Study population: 22 primary care providers in Vietnam enrolled in the StITCH hepatitis care training program; 132 vignette interactions across ten clinical scenarios (five general conditions, five hepatitis-related).

\newpage

\begin{table}[htbp]
\centering
\begin{tabularx}{\textwidth}{|p{0.8cm}|p{5.5cm}|p{6.5cm}|p{3cm}|}
\hline
\textbf{Item} & \textbf{STARD Item} & \textbf{Reported on Page/Section} & \textbf{Notes} \\
\hline

\multicolumn{4}{l}{\textbf{TITLE AND ABSTRACT}} \\
\hline

1 & Identification as a diagnostic accuracy study in the title or abstract & Title, Abstract & Title explicitly identifies as validation study \\
\hline

\multicolumn{4}{l}{\textbf{INTRODUCTION}} \\
\hline

2 & At least one research question or aim stated in the Introduction & Introduction, p.\ 2 & Research question clearly stated: LLM coding reliability vs.\ human-coded ground truth \\
\hline

\multicolumn{4}{l}{\textbf{METHODS}} \\
\hline

3 & Explicit description of study population (inclusion/exclusion criteria) & Methods: Study design, p.\ 3 & 22 providers enrolled in StITCH program; focus group vs.\ validation cohorts clearly distinguished \\
\hline

4 & Index test fully described; sufficient detail for replication & Methods: Data extraction and coding, p.\ 4 & Model specified (claude-haiku-4-5); coding schema described; prompts available in Supporting Information \\
\hline

5 & Reference standard fully described; sufficient detail for replication & Methods: Data extraction and coding, p.\ 4 & Manual coding schema explicitly specified; identical to index test classification logic \\
\hline

6 & Rationale for choice of reference standard & Methods: Data extraction and coding, p.\ 4 & Human coding established as ground truth for validation; manual coding by research team members \\
\hline

7 & Definition of and rationale for any algorithmic steps (e.g., combination of tests) & Methods: Outcomes, p.\ 4 & Checklist score calculated as proportion of items completed; both human-coded and LLM-coded measures reported separately \\
\hline

8 & Blinding: whether assessors of index test knew results of reference standard, and vice versa & Methods: Study design, p.\ 3 & Not applicable; LLM coding and human coding are independent processes by design; no information about order of coding provided \\
\hline

9 & Timing: was there an appropriate interval between index test and reference standard? & Methods: Data extraction and coding, p.\ 4 & LLM coding performed after human coding; separation ensures independence of assessments \\
\hline

10 & Study population: appropriate spectrum of participants to address research question & Methods: Study design, p.\ 3 & Validation sample: 22 providers of varying experience; focus group: 9 providers for usability testing \\
\hline

11 & Study population: how were participants selected? & Methods: Study design, p.\ 3 & Validation cohort: providers from StITCH hepatitis training program; focus group: recruited from provincial hospital \\
\hline

12 & Study population: were inclusion/exclusion criteria clearly described? & Methods: Study design, p.\ 3 & Inclusion: primary care providers; focus group vs.\ validation cohorts clearly separated \\
\hline

13 & Participant flow: number and characteristics of participants who were enrolled & Results: Validation sample, p.\ 5 & 22 providers in validation cohort; 9 in focus group; all completed assigned interactions \\
\hline

14 & Participant flow: number and characteristics of participants for whom the index test was available & Results: Validation sample, p.\ 5 & All 132 interactions were coded by both LLM and human; no missing data \\
\hline

15 & Participant flow: number of participants for whom the reference standard was available & Results: Validation sample, p.\ 5 & All 132 interactions were manually coded; 100\% coverage; no exclusions \\
\hline

16 & Participant flow: number and characteristics of dropouts and those with uninterpretable results & Results: Validation sample, p.\ 5 & No dropouts or exclusions; all 132 interactions retained; no uninterpretable results \\
\hline

\multicolumn{4}{l}{\textbf{ANALYSIS AND RESULTS}} \\
\hline

17 & Calculation of sample size (with power) or explanation if not performed & Methods: Study design, p.\ 3 & Sample size justified by data quantity: 22 providers $\times$ 6 cases = 132 interactions; 433 unique data points; statistical hypothesis testing not employed \\
\hline

18 & Handling of missing data: describe any imputations and exclusions & Methods: Data extraction and coding, p.\ 4 & No missing data; all 132 interactions retained for analysis \\
\hline

19 & Handling of uninterpretable test results: describe how handled & Results: Validation sample, p.\ 5 & No uninterpretable results; all interactions yielded codable transcripts \\
\hline

20 & Statistical measures of diagnostic accuracy (sensitivity, specificity, predictive values) & Results: Validation sample, p.\ 6 & Item-level sensitivity and specificity reported (Figure 3); mean sensitivity 0.31 (95\% CI: 0.26--0.35); mean specificity 0.85 (95\% CI: 0.83--0.87) \\
\hline

21 & Measures of statistical uncertainty (confidence intervals) & Results: Validation sample, p.\ 5--6 & 95\% confidence intervals provided for all point estimates; OLS $r^{2}$ and correlation coefficients reported \\
\hline

22 & Assessment of agreement between index test and reference standard & Results: Validation sample, p.\ 5--6 & Pearson correlation coefficient ($\rho$) and $r^{2}$ reported; item-level sensitivity/specificity plotted; proportion of agreement by item also assessed \\
\hline

23 & Link between index test results and clinical decision-making (if applicable) & Discussion, p.\ 7 & Automated coding could enable scalable competency assessment without human review (though low sensitivity noted as limitation) \\
\hline

\multicolumn{4}{l}{\textbf{DISCUSSION}} \\
\hline

24 & Summary of main findings & Discussion, p.\ 7 & LLM coding correlated moderately with human coding ($\rho$ = 0.51--0.53); directly from Vietnamese comparable to translated English \\
\hline

25 & Discussion of limitations, including sources of potential bias & Discussion: Limitations, p.\ 7--8 & Six major limitations described: small sample size, single-country setting, low mean sensitivity (0.31), lack of direct observation, text-only format, model version variation \\
\hline

26 & Discussion of implications for clinical practice & Discussion, p.\ 7 & LLM-driven vignettes enable low-cost, scalable competency assessment; applicable to diverse health system contexts \\
\hline

27 & Comparison with other diagnostic studies & Discussion, p.\ 7 & Results consistent with recent benchmarks of LLM-based clinical scoring; moderate agreement typical of free-text response systems \\
\hline

28 & Consistency of results across subgroups (if applicable) & Results: Validation sample, p.\ 6 & Item-level analysis shows clustering by sensitivity/specificity patterns; comparison of English vs.\ Vietnamese performance in Figure 3 \\
\hline

\multicolumn{4}{l}{\textbf{SUMMARY AND CONCLUSIONS}} \\
\hline

29 & How results are consistent with other diagnostic studies & Discussion, p.\ 7 & Findings align with literature on LLM-based automated scoring; moderate correlation typical; low sensitivity reflects nuanced interpretation challenges \\
\hline

30 & Reporting of estimates of diagnostic accuracy separately for reader interpretation (e.g., in readers with different experience) & Results: Validation sample, p.\ 5 & Results stratified by comparison (human vs.\ expert ratings; English vs.\ Vietnamese LLM coding); separate figures provided for each \\
\hline

\end{tabularx}
\end{table}

\newpage

\section*{STARD Applicability and Modifications}

This study applies the STARD guidelines to the LLM-versus-human coding validation, which constitutes a diagnostic accuracy assessment. However, the broader study design includes additional components (focus group usability testing, expert clinician ratings) that extend beyond the standard STARD framework. Key modifications and notes:

\begin{itemize}

\item \textbf{Blinding (Item 8):} Blinding is not applicable in the traditional sense, as LLM coding and human coding represent independent processes rather than sequential assessments by different evaluators of the same test. The system was designed so that both processes could operate independently without one informing the other.

\item \textbf{Index Test and Reference Standard (Items 4--5):} The index test (LLM automated coding) and reference standard (human coding) use identical classification schemas, ensuring comparability. Both are applied to the same artifacts (transcripts) rather than to distinct clinical assessments.

\item \textbf{Spectrum of Participants (Item 10):} Participants were all primary care providers from a specific training program (StITCH), which represents a relatively homogeneous population. This limits generalizability but provides a well-characterized sample appropriate for a pilot validation study.

\item \textbf{Sample Size (Item 17):} No a priori power calculation was performed. The sample size was justified by pragmatic considerations (provider recruitment within the StITCH program) and the resulting data quantity (132 interactions, 433 checklist items). The study was framed as a pilot and validation study without formal statistical hypothesis testing.

\item \textbf{Expert Validation (beyond STARD):} The study also compared human-coded checklist scores against expert clinician ratings as an external validation criterion (Pearson's $\rho$ = 0.49--0.59), providing evidence that the underlying human-coded measure is clinically meaningful, not just reliable.

\end{itemize}

\newpage

\section*{Key Diagnostic Accuracy Results}

\subsection*{Item-Level Agreement (Sensitivity and Specificity)}

\begin{itemize}

\item \textbf{Mean sensitivity:} 0.31 (95\% CI: 0.26--0.35) -- the LLM correctly identified approximately 31\% of checklist items that providers actually completed

\item \textbf{Mean specificity:} 0.85 (95\% CI: 0.83--0.87) -- the LLM correctly identified approximately 85\% of checklist items that providers did not complete

\item \textbf{Interpretation:} The LLM shows high specificity but low sensitivity, indicating a tendency to under-call completed items. This would systematically underestimate provider competency if used without human review.

\end{itemize}

\subsection*{Observation-Level Correlation (Checklist Score Agreement)}

\begin{itemize}

\item \textbf{English transcripts:} Pearson's $\rho$ = 0.53; OLS $r^{2}$ = 0.28 -- moderate correlation

\item \textbf{Vietnamese transcripts:} Pearson's $\rho$ = 0.51; OLS $r^{2}$ = 0.26 -- comparable performance without translation

\item \textbf{Interpretation:} A 10 percentage point increase in human-coded checklist completion predicted a 5.9 p.p.\ increase in LLM-coded completion, indicating modest but meaningful agreement at the interaction level.

\end{itemize}

\subsection*{External Validation: Expert Clinician Ratings}

\begin{itemize}

\item \textbf{Human checklist vs.\ expert rating (Vietnamese):} Pearson's $\rho$ = 0.49

\item \textbf{Human checklist vs.\ expert rating (English):} Pearson's $\rho$ = 0.59

\item \textbf{Two expert raters:} Pearson's $\rho$ = 0.62

\item \textbf{Interpretation:} Human-coded checklists correlate moderately to well with independent expert assessments, suggesting the underlying measure captures meaningful clinical performance variation.

\end{itemize}

\end{document}
